\section{Refraction, absorption, and reflectivity}
\label{sec:si_optics}

The main-text reflectivity discussion separates the low-angle optical response from finite-stack diffraction. The equations below specify the propagation factors and the practical crossover used in the physical diagnostic.

\subsection{Wave propagation, refraction, and absorption}
\label{subsec:si_refraction_absorption}

Let $\zeta\geq0$ denote normal depth into the film and use the field convention
$\exp[i(\widetilde{\mathbf k}\cdot\mathbf r-\omega_{\mathrm{em}}t_{\mathrm{time}})]$. For wavelength sample
$\lambda_s$ with source weight $w_s$,
\begin{equation}
  k_s=\frac{2\pi}{\lambda_s},
  \qquad
  \tilde n_q(\lambda_s)=1-\delta_q(\lambda_s)+i\beta_q(\lambda_s),
  \label{eq:si_complex_n_lambda}
\end{equation}
Here $\widetilde{\mathbf k}$ is the complex optical wavevector, $\mathbf r$ is position and $\tilde n_q$ is the complex refractive index. The prime on an optical wavevector identifies propagation inside the film, as in the main text. The index $s$ is the source-state index, $q$ identifies the optical layer, $\omega_{\mathrm{em}}$ is angular frequency, and $t_{\mathrm{time}}$ is time. The real internal phase vectors used in the main text and Sec.~\ref{sec:si_detector_profiles} are $\mathbf k'_{i,s}=\operatorname{Re}\widetilde{\mathbf k}'_{i,s}$ and $\mathbf k'_{f,s}=\operatorname{Re}\widetilde{\mathbf k}'_{f,s}$. The symbol $\widetilde k_\perp$ below denotes the complex component along the local inward normal. The laboratory $z$ axis is defined separately in Sec.~\ref{sec:si_geometric_ray_acceptance}. The decrement $\delta_q$ controls refraction
and critical-angle behavior, whereas $\beta_q$ controls absorption and
penetration depth. Both are wavelength-dependent material inputs. Their
source, interpolation, and association with the sampled spectrum are recorded
with the calculation configuration.

\paragraph{Tangential matching and the complex normal component.}
Translational invariance along a flat interface conserves the tangential
wavevector. For a field entering layer $q$ from the ambient medium (index $0$), let $\mathbf k_{\parallel,q}$ denote the conserved tangential component. The superscript $(+)$ selects the inward, decaying normal root:
\begin{equation}
  \mathbf k_{\parallel,q}(\lambda_s)
  =\mathbf k_{\parallel,0}(\lambda_s),
  \qquad
  \widetilde k_{\perp,q}^{(+)}(\lambda_s)
  =
  \sqrt{
  \left[\tilde n_q(\lambda_s)k_s\right]^2
  -\left|\mathbf k_{\parallel,0}(\lambda_s)\right|^2}.
  \label{eq:si_boundary_ktz_lambda}
\end{equation}
The root is chosen so that the field decays in its propagation direction. For
the inward solution and the depth convention above,
$\operatorname{Im}\widetilde k_{\perp,q}^{(+)}\geq0$.\cite{Renaud2009} The reciprocal exit
field is obtained from the same boundary condition with its own conserved
tangential component and propagation direction. Real transmitted angles may be reported after this complex-vector calculation. Their evaluation retains both the real and imaginary parts of $\tilde n_q$.

\paragraph{Interface amplitudes and normal flux.}
For a scalar field at an interface from medium $p$ to medium $q$, the field
transmission amplitude is
\begin{equation}
  t_{pq}=\frac{2\widetilde k_{\perp,p}}{\widetilde k_{\perp,p}+\widetilde k_{\perp,q}}.
  \label{eq:si_scalar_transmission}
\end{equation}
Equation~\eqref{eq:si_scalar_transmission} specifies the field-amplitude coefficient. Converting it to a transmitted-flux fraction requires the appropriate flux normalization. In the nonabsorbing scalar limit, the associated
normal-flux ratio includes the factor
$\operatorname{Re}\widetilde k_{\perp,q}/\operatorname{Re}\widetilde k_{\perp,p}$ in addition to
$|t_{pq}|^2$. Absorbing or polarization-resolved layers require the
corresponding Poynting-flux or Fresnel normalization. The calculation must
therefore declare whether each interface weight is a field amplitude, an
intensity, or a normal flux. The transmitted-flux fraction combines $|t_{pq}|^2$ with that normalization.

\paragraph{Normal-depth attenuation.}
Let $\kappa_i(\lambda_s)$ and $\kappa_f(\lambda_s)$ be the nonnegative
imaginary normal components for propagation from the surface to a scattering
depth and from that depth back to the surface, respectively:
\begin{equation}
  \kappa_i(\lambda_s)=\operatorname{Im}\widetilde k_{\perp,i}^{(+)}(\lambda_s),
  \qquad
  \kappa_f(\lambda_s)=\operatorname{Im}\widetilde k_{\perp,f}^{(+)}(\lambda_s),
  \qquad
  \kappa_i,\kappa_f\geq0.
  \label{eq:si_kappa_if_lambda}
\end{equation}
Here \(\widetilde k_{\perp,f}^{(+)}\) is the inward reciprocal exit-channel root used by
reciprocity. Thus \(\kappa_f\) is the positive attenuation constant for the
physical path from depth back to the surface. Equivalently,
\(\kappa_f=-\operatorname{Im}\widetilde k_{\perp,f}^{(\mathrm{out})}\) if an outward
wavevector is stored explicitly.
For a normalized scattering-depth distribution $\rho_{\mathrm{depth}}(\zeta)$ over a film of
thickness $t$, the propagation weight is
\begin{equation}
  W_{\mathrm{depth}}(\lambda_s)=
  \int_0^t \rho_{\mathrm{depth}}(\zeta)
  \exp\!\left[-2\left(\kappa_i(\lambda_s)+\kappa_f(\lambda_s)\right)\zeta\right]d\zeta,
  \qquad
  \int_0^t\rho_{\mathrm{depth}}(\zeta)\,d\zeta=1.
  \label{eq:si_depth_weight_general}
\end{equation}
For uniform scattering density, this becomes
\begin{align}
  W_{\mathrm{depth}}(\lambda_s)
  &=
  \frac{1}{t}\int_0^t
  \exp\!\left[-2\left(\kappa_i(\lambda_s)+\kappa_f(\lambda_s)\right)\zeta\right] d\zeta\\
  &=
  \frac{
  1-\exp\!\left[-2\left(\kappa_i(\lambda_s)+\kappa_f(\lambda_s)\right)t\right]
  }
  {2\left(\kappa_i(\lambda_s)+\kappa_f(\lambda_s)\right)t}.
  \label{eq:si_abs_complex_kz}
\end{align}
Equations~\eqref{eq:si_depth_weight_general} and
\eqref{eq:si_abs_complex_kz} use normal depth. Multiplying
$\operatorname{Im}\widetilde k_\perp$ by an oblique ray-path length would count the angular
projection a second time. The depth weight is combined with consistently
normalized entrance and exit interface factors before it multiplies the
structure-factor and mosaic weights.

\paragraph{Detector accumulation over the bandwidth.}
For each accepted source-state and rod contribution, the wavelength-resolved detector density is formed with
the same geometry, structure-factor, mosaic, detector, and optical interfaces.
The detector images are then added incoherently,
\begin{equation}
  I_{\mathrm{det}}(c_D,r_D)
  =
  \sum_s
  w_s\,
  I_{\mathrm{det},s}\!\left[c_D,r_D;\lambda_s,\{\tilde n_q(\lambda_s)\}\right].
  \label{eq:si_detector_bandwidth_sum}
\end{equation}
Changing $\lambda_s$ changes the moving Ewald construction, the layer optical
constants, the complex normal components, and the interface and depth weights.
The planar-interface correction is retained for every wavelength state. Its
role is expected to be strongest when an incident or exit angle approaches the
critical-angle regime. This optical treatment supplies a propagation correction. Assessing experimental intensity agreement requires measured/calculated comparisons.


\subsection{Parratt-to-kinematic crossover for the axial profile}
\label{subsec:si_parratt_structure_crossover}

The detector-derived $r=0$ profile spans two limiting regimes. The low-$Q_{\mathrm{loc}}$ response is described by standard Parratt reflectivity, while the higher-$Q_{\mathrm{loc}}$ Bragg response is described by a finite-stack kinematic term. The two branches are matched smoothly over a declared overlap interval. We refer to this matching as the Parratt-to-kinematic crossover. It is a practical connection between two established limiting calculations. The main-text reflectivity figure shows measured Bi$_2$Se$_3$ intensities versus detector $2\theta$.

For the physical crossover diagnostic, the dimensionless coordinate is
\begin{equation}
  x=\frac{Q_{\mathrm{loc}}}{Q_c},
  \qquad
  L_{\mathrm{loc}}=\frac{Q_{\mathrm{loc}} c}{2\pi}=x\frac{Q_c c}{2\pi}.
  \label{eq:si_crossover_x_l}
\end{equation}
Here $Q_c$ is the film critical momentum transfer for total external reflection at wavelength $\lambda_s$, and $c$ is the out-of-plane lattice
constant. The transfer $Q_{\mathrm{loc}}=|\mathbf k_{f,s}-\mathbf k_{i,s}|$ is the local momentum-transfer magnitude for specular reflection from a crystallite, whose normal is parallel to that transfer. It generally differs from its mean-film projection $Q_z$. At nominal incidence $\theta_i$, the mean-film specular direction is $2\theta=2\theta_i$. Tilted crystallites can contribute axial scattering away from that direction. The $L_{\mathrm{loc}}$ value in Eq.~\eqref{eq:si_crossover_x_l} is an external local optical coordinate. Finite-stack scattering uses the separate intrinsic phase coordinate. The diagnostic crossover coordinate uses $Q_{\mathrm{loc}}/Q_c$ and the
corresponding $L_{\mathrm{loc}}$ value defined above.

The low-$Q_{\mathrm{loc}}$ branch is the Parratt reflectivity.\cite{Parratt1954} For a layer
$q$ with complex refractive index $\tilde n_q=1-\delta_q+i\beta_q$, the conserved
in-plane wavevector fixes the layer-normal component
\begin{equation}
  \widetilde k_{\perp,q}=\sqrt{(\tilde n_q k_s)^2-\left|\mathbf k_{\parallel}\right|^2},
  \qquad
  k_s=\frac{2\pi}{\lambda_s}.
  \label{eq:si_parratt_kz}
\end{equation}
The fixed source-state label $s$ is suppressed on the layer amplitudes below. For the depth and field convention used above, the square-root branch is chosen
to give decay in the relevant propagation direction.
The superscript $(0)$ labels the smooth-interface coefficient. The Fresnel amplitude at the interface between layers $q$ and $q+1$ is
\begin{equation}
  r_q^{(0)}=\frac{\widetilde k_{\perp,q}-\widetilde k_{\perp,q+1}}{\widetilde k_{\perp,q}+\widetilde k_{\perp,q+1}}.
  \label{eq:si_parratt_interface}
\end{equation}
If the lower interface has RMS roughness $\sigma_{\mathrm{rough},q}$, this coefficient is damped
by
\begin{equation}
  r_q=r_q^{(0)}\exp\!\left[-2\widetilde k_{\perp,q}\widetilde k_{\perp,q+1}\sigma_{\mathrm{rough},q}^2\right].
  \label{eq:si_parratt_roughness}
\end{equation}
For a finite layer of thickness $d_{q+1}$ below the interface, the round-trip
phase factor $P_{\mathrm{opt},q+1}$ is
\begin{equation}
  P_{\mathrm{opt},q+1}=\exp\!\left(2i\widetilde k_{\perp,q+1}d_{q+1}\right).
  \label{eq:si_parratt_phase}
\end{equation}
Starting from the semi-infinite substrate and recursing upward, the effective
reflection amplitude $R_q^{\mathrm{amp}}$ at interface $q$ is
\begin{equation}
  R_q^{\mathrm{amp}}
  =
  \frac{r_q+R_{q+1}^{\mathrm{amp}}P_{\mathrm{opt},q+1}}
       {1+r_qR_{q+1}^{\mathrm{amp}}P_{\mathrm{opt},q+1}}.
  \label{eq:si_parratt_recursion}
\end{equation}
The optical branch in this diagnostic is
\begin{equation}
  R_{\mathrm P}(Q_{\mathrm{loc}})=\left|R_0^{\mathrm{amp}}(Q_{\mathrm{loc}})\right|^2.
  \label{eq:si_parratt_reflectivity}
\end{equation}
This branch retains total external reflection, refraction, evanescence, finite
penetration depth, roughness damping, and multiple reflection between the film
interfaces.

The high-$Q_{\mathrm{loc}}$ branch is built from the normalized finite-stack kinematic term
\begin{equation}
  \widehat S_{00}(L)=
  \frac{S_{q,00}(|\mathbf c^*|L)}{S_{q,00}(0)}.
  \label{eq:si_ht_normalized_structure}
\end{equation}
The hat denotes normalization to the zero-transfer strength, so $\widehat S_{00}(0)=1$. The fixed wavelength $\lambda_s$ is suppressed in these axial structure terms. Here $S_{q,00}$ is the ordered axial structural strength per dimensional rod coordinate from Sec.~\ref{subsec:si_rod_measure}. Near the critical edge, the structure term uses the internal film phase. To match the Parratt film phase, the finite-stack coordinate is computed from the real propagating part of the film wavevector,
\begin{equation}
  Q_{\mathrm{film}}=\operatorname{Re}\!\left(2\widetilde k_{\perp,\mathrm{film}}\right),
  \qquad
  L_{\mathrm{film}}=\frac{Q_{\mathrm{film}}c}{2\pi}.
  \label{eq:si_internal_phase_coordinate}
\end{equation}
Below the critical edge this real part may approach zero in the weak-absorption
limit, while a complex internal wavevector still describes evanescence and
absorption. The superscript $(0)$ on the kinematic branch below labels its unscaled value:
\begin{equation}
  R_{\mathrm{kin}}^{(0)}(Q_{\mathrm{loc}})=\frac{\widehat S_{00}(L_{\mathrm{film}})}{Q_{\mathrm{loc}}^2}.
  \label{eq:si_crossover_unscaled_structure}
\end{equation}
Since $\widehat S_{00}$ is dimensionless, $R_{\mathrm{kin}}^{(0)}$ has units of length squared. It is an unscaled intensity shape. The $Q_{\mathrm{loc}}^{-2}$ factor gives the high-$Q_{\mathrm{loc}}$ structural envelope, while the
internal coordinate $L_{\mathrm{film}}$ keeps the finite-stack fringe phase consistent with
the optical film phase.

A multiplicative scale places this kinematic branch on the Parratt intensity
scale. The scale is fit in a high-$Q_{\mathrm{loc}}$ overlap window, excluding the selected
$00L$ neighborhoods used in the diagnostic. With fit-point index $i$ and retained fit set $\mathcal F$, the
log-median estimate is
\begin{equation}
  A_{\mathrm{opt}}=Q_c^2
  \exp\!\left[
  \operatorname{median}_{i\in\mathcal F}
  \ln\!\left(
  \frac{R_{\mathrm P}(Q_{\mathrm{loc},i})}
       {Q_c^2R_{\mathrm{kin}}^{(0)}(Q_{\mathrm{loc},i})}
  \right)
  \right],
  \label{eq:si_crossover_scale}
\end{equation}
with the default overlap window $5<Q_{\mathrm{loc}}/Q_c<10$. The factors of $Q_c^2$ make the logarithm's argument dimensionless and give $A_{\mathrm{opt}}$ units of inverse length squared. The scaled kinematic branch is therefore dimensionless, like $R_{\mathrm P}$:
\begin{equation}
  R_{\mathrm{kin}}(Q_{\mathrm{loc}})=A_{\mathrm{opt}} R_{\mathrm{kin}}^{(0)}(Q_{\mathrm{loc}}).
  \label{eq:si_crossover_scaled_structure}
\end{equation}

The crossover window is selected by comparing the two positive finite branches in
log intensity. Both branches are evaluated at the same dimensional transfer $Q_{\mathrm{loc}}=Q_c x$:
\begin{equation}
  \delta_{\log}(x)=\log_{10}\!\left[\frac{R_{\mathrm{kin}}(Q_c x)}{R_{\mathrm P}(Q_c x)}\right].
  \label{eq:si_crossover_log_ratio}
\end{equation}
Candidate points must lie in the search range $3\le x\le10$ and satisfy
$|\delta_{\log}(x)|\le0.10$. The selected window $[x_1,x_2]$ is the widest contiguous
eligible region with width at least $1.0$ in $Q_{\mathrm{loc}}/Q_c$. Ties are broken by the
smaller median $|\delta_{\log}|$. If the search rejects every candidate region, the fallback window
is $x_1=3$ and $x_2=6$.

Inside the selected window, the interpolation coordinate and weight are
\begin{equation}
  \eta_{\mathrm{blend}}=\frac{x-x_1}{x_2-x_1},
  \qquad
  w_{\mathrm{blend}}(\eta_{\mathrm{blend}})=6\eta_{\mathrm{blend}}^5-15\eta_{\mathrm{blend}}^4+10\eta_{\mathrm{blend}}^3,
  \label{eq:si_crossover_smoothstep}
\end{equation}
with $\eta_{\mathrm{blend}}$ clipped to the interval $[0,1]$. The final branch is blended in log
reflectivity,
\begin{equation}
R_{\mathrm{blend}}(x)=
\begin{cases}
R_{\mathrm P}(Q_c x), & x\le x_1,\\[3pt]
10^{(1-w_{\mathrm{blend}})\log_{10}R_{\mathrm P}(Q_c x)+w_{\mathrm{blend}}\log_{10}R_{\mathrm{kin}}(Q_c x)},
  & x_1<x<x_2,\\[3pt]
R_{\mathrm{kin}}(Q_c x), & x\ge x_2.
\end{cases}
\label{eq:si_crossover_blend}
\end{equation}
A small positive numerical floor is used only to avoid logarithms of nonpositive
roundoff values inside the blend window. The resulting dimensionless reflectivity is exactly
Parratt below the selected window, exactly the scaled structure branch above the
window, and smooth through the crossover.


The subscript $\mathrm{stitch}$ labels the replacement axial strength after this crossover. The dimensionless blend is converted to the axial structural strength before detector weighting. In the per-$dq_z$ convention of the optical calculation,
\begin{equation}
S_{q,\mathrm{stitch}}(Q_{\mathrm{loc}})
=\frac{Q_{\mathrm{loc}}^2 S_{q,00}(0)}{A_{\mathrm{opt}}}R_{\mathrm{blend}}(Q_{\mathrm{loc}}/Q_c).
\label{eq:si_reflectivity_strength_conversion}
\end{equation}
Above the crossover this recovers the finite-stack strength evaluated at $L_{\mathrm{film}}$. The stitch supplies the replacement axial structural strength and enters the detector weighting once. Equation~\eqref{eq:si_rod_measure_conversion} supplies the conversion to the per-$dL$ density used in Sec.~\ref{sec:si_detector_profiles}.

For the joint comparison, the low-$Q_{\mathrm{loc}}$ interval and the $003$ neighborhood share one normalization. The measured detector profile, the Parratt
branch, the finite-stack branch, and the combined curve must share one declared
background treatment and one intensity scale. The $003$ peak serves as the
higher-$Q_{\mathrm{loc}}$ structural reference. Its sensitivity to wavelength spread can be examined
with Eqs.~\eqref{eq:si_full_spectrum_control} and
\eqref{eq:si_mean_wavelength_control}. The model controls described above
retain the radial extension for a monochromatic source. The contribution of wavelength broadening to the measured feature remains to be established.
Figure~\ref{fig:si_optical_comparison} reserves the physical comparison on that common scale.


\begin{figure}[!htbp]
  \centering
  \siplaceholder{Physical axial reflectivity comparison}{Measured axial profile, direct Parratt and finite-stack branches, combined prediction, and residual.\\Common coordinate map, background, and intensity scale.}
  \caption{Placeholder for the physical axial comparison. The low-angle interval and diffraction region are to be calculated from one declared structural and optical state with uncertainty analysis.}
  \label{fig:si_optical_comparison}
\end{figure}
